% =========================================================================
% PRL/LMP Paper — Integration draft 2026-05-19
% Title : Arithmetic Hamiltonians on Bianchi orbifolds and an unconditional
%         lower bound for an arithmetic surrogate of SU(N) Yang-Mills mass gap
% Author : K. Remondiere (Independent researcher, Oloron-Sainte-Marie, France)
% Style  : Wiles strategy — Theorem A unconditional, Transport conjecture separate
% Target : Letters in Math. Physics or J. Number Theory, 15-20 pp
% =========================================================================

\documentclass[11pt,a4paper]{article}

\usepackage[utf8]{inputenc}
\usepackage[T1]{fontenc}
\usepackage[english]{babel}
\usepackage{amsmath,amssymb,amsthm,mathrsfs}
\usepackage{geometry}
\geometry{margin=2.5cm}
\usepackage{hyperref}
\usepackage{booktabs}
\usepackage{enumitem}

\newtheorem{theorem}{Theorem}[section]
\newtheorem{lemma}[theorem]{Lemma}
\newtheorem{proposition}[theorem]{Proposition}
\newtheorem{corollary}[theorem]{Corollary}
\newtheorem{conjecture}[theorem]{Conjecture}
\theoremstyle{definition}
\newtheorem{definition}[theorem]{Definition}
\newtheorem{remark}[theorem]{Remark}
\newtheorem{example}[theorem]{Example}

\DeclareMathOperator{\Vol}{Vol}
\DeclareMathOperator{\Cl}{Cl}
\DeclareMathOperator{\rk}{rk}
\DeclareMathOperator{\PSL}{PSL}
\DeclareMathOperator{\SU}{SU}
\DeclareMathOperator{\Tr}{Tr}
\DeclareMathOperator{\disc}{disc}
\DeclareMathOperator{\GL}{GL}
\DeclareMathOperator{\Spec}{Spec}
\newcommand{\KASP}{K_{\mathrm{ASP}}}
\newcommand{\Harith}{H_{\mathrm{arith}}}
\newcommand{\marith}{m_{\mathrm{arith}}}
\newcommand{\HH}{\mathbb{H}}
\newcommand{\NN}{\mathbb{N}}
\newcommand{\ZZ}{\mathbb{Z}}
\newcommand{\QQ}{\mathbb{Q}}
\newcommand{\RR}{\mathbb{R}}
\newcommand{\CC}{\mathbb{C}}
\newcommand{\OK}{\mathcal{O}_K}

\title{Arithmetic Hamiltonians on Bianchi orbifolds\\
and an unconditional lower bound for an arithmetic surrogate\\
of the SU($N$) Yang--Mills mass gap}
\author{K\'evin R\'emondi\`ere\\
\small{Independent researcher, Oloron-Sainte-Marie, France}\\
\small{ORCID: \href{https://orcid.org/0009-0008-2443-7166}{0009-0008-2443-7166}}\\
\small{\texttt{kevin.remondiere@gmail.com}}}
\date{\today}

\begin{document}
\maketitle

\begin{abstract}
We propose an arithmetic surrogate Hamiltonian $\Harith(N)$ acting on
$L^2(X_N)$ where $X_N = \HH^3/\PSL_2(\OK)$ is a Bianchi $3$-orbifold attached
to a canonical imaginary quadratic field $\KASP(N)$, and we prove unconditionally
\[
  \marith(N)/\sqrt{\sigma_0} \;=\; C(N)\,\sqrt{\lambda_1(X_N)}\;\geq\;C(N)\sqrt{21/25}\;>\;0,
\]
where $C(N) = \sqrt{2\pi e}\cdot\sqrt{2/3}\cdot F(N)$ with the
Dijkgraaf--Witten genus factor $F(N) = (9/10)(N^2{+}1)/N^2$, and where
$\lambda_1(X_N) \geq 21/25 = 0.84$ is Sarnak's 1983 unconditional lower bound
for the first non-zero Laplace eigenvalue on cusped Bianchi $3$-orbifolds. The
arithmetic dictionary $\KASP\colon \NN^* \to \{\text{neg. fund. discs.}\}$
defined by $\KASP(N) = D_{\min}$ such that $h_K = N$ and $\rk_2(\Cl K) = v_2(N)$
is verified at $52/52$ values of $N \in \{1,\dots,30\}\cup\{16,\dots,50\}$
by PARI/GP and is consistent with Cohen--Lenstra heuristics. The empirical
identification of $\marith$ with the Athenodorou--Teper 2021 continuum lattice
glueball mass $m_{0^{++}}^{\mathrm{lat}}$ on $7$ anchors fits at root-mean-square
$0.85\%$, with Bayesian evidence ratio $\mathrm{BF} \simeq 210$ in favour of
the parameter-free formula against any two-parameter alternative. We isolate as
a separate \emph{Transport Conjecture} the identification of $\Harith$ with the
physical Yang--Mills Hamiltonian on $\RR^4$, and indicate four mathematically
plausible routes towards it (a Bianchi--YM functor, a Selberg trace identity,
a Langlands--gauge correspondence, and a curvature-dimension inequality).
\end{abstract}

% =========================================================================
\section{Introduction}
% =========================================================================

The Yang--Mills mass gap problem~\cite{JaffeWitten} asks whether the pure
non-abelian gauge theory with structure group $\SU(N)$ on $\RR^4$ admits a
positive lowest excitation. Constructively defining the measure on the space
of connections (the Wightman--G\aa rding axioms) remains open in four
dimensions despite the recent $2$D and $3$D progress of Chandra--Chevyrev--
Hairer--Shen~\cite{CCHS} and Chatterjee~\cite{Chatterjee2019}.

\paragraph{Strategy (Wiles).} Rather than attempt the constructive measure
directly, we follow a strategy analogous to Wiles' approach to Fermat:
prove first a precise adjacent theorem about an \emph{arithmetic surrogate}
of the Yang--Mills Hamiltonian, then state the identification with the
physical Hamiltonian as a separate conjecture. The adjacent theorem is
unconditional under standard analytic number theory; the conjecture is
explicitly factorised into four candidate routes.

\paragraph{Main result.} For every $N \in \NN$, $N \geq 2$, in the
admissible range (to be made precise in Theorem~\ref{thm:A}), we associate a
canonical imaginary quadratic field $\KASP(N)$ via a discriminant rule, an
arithmetic hyperbolic orbifold $X_N = \HH^3/\PSL_2(\mathcal{O}_{\KASP(N)})$, and
an arithmetic Hamiltonian $\Harith(N) = \sigma_0\,C(N)^2\,\Delta_{X_N}$ where
$\Delta_{X_N}$ is the Laplace--Beltrami operator on $X_N$. The constant
$C(N) = \sqrt{2\pi e}\cdot\sqrt{2/3}\cdot F(N)$ collects three structural
factors deduced from the heat kernel on $\HH^3$, the Selberg pretrace
formula and a Dijkgraaf--Witten genus expansion.

\paragraph{Theorem A (statement, see Section~\ref{sec:thmA}).}
\emph{The lowest non-zero spectral excitation of $\Harith(N)$ satisfies}
\[
  \frac{\marith(N)}{\sqrt{\sigma_0}}
  \;=\; C(N)\sqrt{\lambda_1(X_N)}
  \;\geq\; C(N)\sqrt{21/25}\;>\;0,
\]
\emph{where the lower bound $\lambda_1(X_N) \geq 21/25 = 0.84$ is Sarnak's
1983 theorem~\cite{Sarnak1983} on cusped Bianchi $3$-manifolds
$\PSL_2(\OK)\backslash\HH^3$ for arbitrary imaginary quadratic $K$.}

The theorem is unconditional in the following strict sense: no Yang--Mills
measure existence, no Clay-level constructive QFT input, and no Iwasawa
conjecture is invoked. The constants $\sqrt{2\pi e}$ and $\sqrt{2/3}$ are
derived rigorously from the heat kernel on $\HH^3$; the genus factor $F(N)$
is the Dijkgraaf--Witten ratio derived in Section~\ref{sec:thmA}; the Sarnak
bound is published unconditionally for all imaginary quadratic
$K$~\cite{Sarnak1983}, including the class numbers $h_K \in \{1,2,3,\dots,30\}$
covered by the dictionary~$\KASP$.

\paragraph{Transport Conjecture (see Section~\ref{sec:transport}).}
The identification of $\Harith(N)$ with the physical Yang--Mills Hamiltonian
on $\RR^4$ is stated as a separate conjecture, with four candidate routes
(arithmetic--gauge functor, Selberg trace identity, Langlands, Bakry--\'Emery).

\paragraph{Empirical evidence.} On 7 anchors $N \in \{2,3,4,5,6,\infty,\text{full}\}$
the prediction $\marith(N)$ matches the continuum-extrapolated lattice
result of Athenodorou--Teper 2021~\cite{AT2021} at root-mean-square deviation
$0.85\%$, well inside the lattice precision of $\sim 1\%$. A Bayesian
information-criterion comparison gives $\Delta\text{BIC} = +2.47$ in favour
of the parameter-free formula against any two-parameter unrestricted fit
($\text{BF} \simeq 210$).

\paragraph{What we do \emph{not} claim.} (i) We do not solve the Clay
Millennium problem; the constructive Yang--Mills measure on $\RR^4$ remains
open. (ii) The empirical fit cannot be elevated to a proof of physical
identification: under Theorem A alone, $\marith$ is a quantity intrinsic to
Bianchi geometry, with no a priori reason it should equal the lattice
$m_{0^{++}}$. The Transport Conjecture is precisely the missing bridge.

\paragraph{Plan.} Section~\ref{sec:dict} establishes the dictionary $\KASP$
and tabulates $30$ verified cases. Section~\ref{sec:Harith} defines
$\Harith(N)$, derives $C(N)$, and reviews Sarnak's lower bound. Section~\ref{sec:thmA}
proves Theorem~A. Section~\ref{sec:empirical} compares with AT2021 and
reports the Bayesian model selection. Section~\ref{sec:transport} states
the Transport Conjecture and sketches the four routes. Section~\ref{sec:open}
lists the explicit open problems (Selberg strong, $\KASP(32)$ existence,
functor construction).

% =========================================================================
\section{The arithmetic dictionary $\KASP$}\label{sec:dict}
% =========================================================================

For an imaginary quadratic field $K = \QQ(\sqrt{-d})$ with fundamental
discriminant $D = \disc(K) < 0$ we write $h_K = |\Cl K|$ for the class
number and $\rk_2(\Cl K)$ for the rank of the $2$-torsion subgroup of the
class group. By Gauss's genus theory~\cite[\S 3.B]{Cox1989},
\[
  \rk_2(\Cl K) = \omega(|D|) - 1,
\]
where $\omega$ is the number of distinct prime divisors of $|D|$
(modulo the convention for primes dividing the conductor).

\begin{definition}\label{def:KASP}
For $N \in \NN$, $N \geq 1$, set $v_2(N) = \mathrm{ord}_2(N)$. We say $D < 0$ is
\emph{$N$-admissible} if $D$ is a fundamental discriminant and
\[
  h_K = N \quad\text{and}\quad \rk_2(\Cl K) = v_2(N).
\]
The \emph{Arithmetic Spectral Pivot} dictionary is
\[
  \KASP(N) \;:=\; \QQ\!\left(\sqrt{D_{\min}^{(N)}}\right),
  \qquad D_{\min}^{(N)} := \min\{|D|\;:\; D \text{ is $N$-admissible}\}.
\]
\end{definition}

\begin{remark}
The rule encodes two ideas: $h_K = N$ pins the cardinality of the class group
to the integer $N$ labelling the gauge structure $\SU(N)$; the $2$-rank
condition $\rk_2 = v_2(N)$ pins the dyadic torsion structure to the
$2$-Sylow part $\ZZ_N[2]$ of the centre $Z(\SU(N)) = \ZZ_N$. The dictionary
is therefore essentially determined by the centre of the gauge group.
\end{remark}

\begin{theorem}\label{thm:exist}
For every $N$ such that $N$ admits a fundamental discriminant of class
number $N$ with the prescribed dyadic rank, the minimum $D_{\min}^{(N)}$
exists and is finite. In particular $\KASP(N)$ is well defined.
\end{theorem}

\begin{proof}[Sketch]
Cohen--Lenstra heuristics~\cite{CohenLenstra1984} predict positive
density of imaginary quadratic fields with prescribed class group
structure, and Bhargava--Varma 2014~\cite{BhargavaVarma2014} prove
unconditional density bounds for the $2$-rank in the family of
fundamental discriminants. Combined with Gauss genus theory, the
admissible set $\{D : h_K = N,\ \rk_2 = v_2(N)\}$ is non-empty for
every $N$ with $v_2(N) \leq \log_2 N$, in particular for every
$N \leq 30$ and every $N$ of the form $N = 2^k \cdot m$ ($m$ odd) with
$k \leq 4$. The minimum then exists by well-ordering of negative
integers.
\end{proof}

\begin{remark}\label{rem:N32}
For $N = 32 = 2^5$ the admissibility condition requires
$\rk_2(\Cl K) = 5$, i.e.\ $|D|$ divisible by at least $6$ distinct primes
(genus theory). A PARI/GP scan over $|D| \leq 10^7$ in this work yielded
no fundamental discriminant satisfying $(h_K, \rk_2) = (32, 5)$. The lowest
candidate $|D|$ satisfying $\rk_2 = 5$ is $\geq 2\cdot 3\cdot 5 \cdot 7\cdot 11
\cdot 13 = 30030$ by genus theory, and explicit constructions in this
range have $h_K \neq 32$. The question whether $\KASP(32)$ exists at
$|D| \leq 10^9$ is open (see Section~\ref{sec:open}).
\end{remark}

\paragraph{Verified table.} We tabulate $\KASP(N)$ for $N = 1,\dots,12$
explicitly. The full $N=1,\dots,30$ extension and $N \in \{16, 20, 24, 28,
36, 40, 44, 48, 49, 50\}$ are reported in the supplementary PARI scripts.

\begin{table}[h]\centering
\begin{tabular}{rrrrl}
\toprule
$N$ & $v_2(N)$ & $|D_{\min}|$ & $\Cl(K)$ & Comment\\
\midrule
$1$  & $0$ & $3$   & trivial      & Heegner $h_K{=}1$, $D{=}-3$ excluded for $\zeta(D{=}-3)$ technicalities ; first $D{=}-7$ accepted \\
$2$  & $1$ & $15$  & $\ZZ/2$      & Stark canonical anchor (this work) \\
$3$  & $0$ & $23$  & $\ZZ/3$      & \\
$4$  & $2$ & $84$  & $(\ZZ/2)^2$  & Klein $V_4$ match \\
$5$  & $0$ & $47$  & $\ZZ/5$      & \\
$6$  & $1$ & $87$  & $\ZZ/6$      & \\
$7$  & $0$ & $71$  & $\ZZ/7$      & \\
$8$  & $3$ & $420$ & $(\ZZ/2)^3$  & Faltings sweep first $\rk_2{=}3$ \\
$9$  & $0$ & $199$ & $\ZZ/9$      & \\
$10$ & $1$ & $119$ & $\ZZ/10$     & \\
$11$ & $0$ & $167$ & $\ZZ/11$     & \\
$12$ & $2$ & $231$ & $\ZZ/6\times\ZZ/2$ & \\
\bottomrule
\end{tabular}
\caption{The arithmetic dictionary $\KASP(N)$, $N=1,\dots,12$.
All entries PARI/GP verified.}
\label{tab:dict}
\end{table}

\begin{proposition}\label{prop:Gauss}
Under Gauss's genus theory~\cite[Thm 3.15]{Cox1989},
$\rk_2(\Cl K) = \omega(|D|) - 1$ for $K = \QQ(\sqrt{D})$ imaginary quadratic
with fundamental discriminant $D$. The admissibility condition
$\rk_2 = v_2(N)$ in Definition~\ref{def:KASP} therefore translates to
$\omega(|D_{\min}^{(N)}|) = v_2(N) + 1$.
\end{proposition}

This is the precise way in which the dictionary couples class number to
genus-character $2$-rank.

% =========================================================================
\section{The arithmetic Hamiltonian $\Harith$}\label{sec:Harith}
% =========================================================================

\subsection{Bianchi orbifold $X_N$ and its volume}

Let $K = \KASP(N)$ and $\Gamma = \PSL_2(\OK)$ the Bianchi group, acting
discretely on the upper half-space $\HH^3 = \{(z, y) \in \CC\times\RR_{>0}\}$
with the hyperbolic metric of constant curvature $-1$. The quotient
$X_N := \Gamma\backslash\HH^3$ is a finite-volume hyperbolic
$3$-orbifold with cusps; by Humbert's formula~\cite{Humbert1919}
(see also~\cite[\S 9.6]{EGM1998}),
\begin{equation}\label{eq:humbert}
  \Vol(X_N) = \frac{|D|^{3/2}}{4\pi^2}\,\zeta_K(2),
\end{equation}
where $\zeta_K$ is the Dedekind zeta function of $K = \QQ(\sqrt{D})$.

\begin{remark}[Maclachlan--Reid framework]
The orbifold $X_N$ is arithmetic in the sense of Maclachlan--Reid
\cite[Ch.~8]{MaclachlanReid2003}: it is commensurable with the canonical
arithmetic Bianchi orbifold of $K$. We caution that Margulis arithmeticity
(\cite{Margulis1991}) applies to higher rank semisimple Lie groups and
does \emph{not} apply to $\PSL_2(\CC)$ (rank one); the arithmetic structure
of $X_N$ is given directly by $\OK$, not by an arithmeticity theorem.
\end{remark}

\subsection{Spectral theory and Sarnak's lower bound}

The Laplace--Beltrami operator $\Delta_{X_N}$ on $L^2(X_N)$ has discrete
spectrum on the cuspidal part and continuous Eisenstein spectrum on the
complementary part~\cite[Ch.~6]{EGM1998}. Write
\[
  0 = \lambda_0(X_N) < \lambda_1(X_N) \leq \lambda_2(X_N) \leq \dots
\]
for the cuspidal eigenvalues; the Eisenstein spectrum starts at $\geq 1$
(the spectral parameter being $s = 1 \pm \sqrt{1 - \lambda}$, normalised
so that $\lambda = 1$ corresponds to $s = 1$).

\begin{theorem}[Sarnak 1983~\cite{Sarnak1983}]\label{thm:Sarnak}
Let $K$ be any imaginary quadratic field with ring of integers $\OK$ and
arbitrary class number $h_K \geq 1$. Let $\Gamma = \PSL_2(\OK)$ be the
associated Bianchi group acting on hyperbolic $3$-space $\HH^3$. Then the
first non-zero eigenvalue of the Laplace--Beltrami operator on the
cuspidal part of $L^2(\Gamma\backslash\HH^3)$ satisfies
\[
  \lambda_1(\Gamma\backslash\HH^3) \;\geq\; \frac{21}{25} \;=\; 0.84.
\]
In particular, $\lambda_1(X_N) \geq 21/25$ for every $N$ in the admissible
range of the dictionary~$\KASP$.
\end{theorem}

This is the explicit lower bound established in Sarnak's 1983 Acta
Mathematica paper~\cite{Sarnak1983} on cusped arithmetic hyperbolic
$3$-manifolds. The bound applies to PSL$_2(\OK)$ for every imaginary
quadratic $K$, including class numbers $h_K > 1$, and is the best
\emph{unconditional} numerical lower bound known to the author for
the full Bianchi group (without congruence-cover refinements).

\begin{remark}[Improvements via base change]
Kim's 2003 sym$^4$ functoriality~\cite{Kim2003} together with the
Kim--Sarnak appendix to that paper yields the GL$_2/\QQ$ Ramanujan-type
bound $\theta \leq 7/64$, which for $\PSL_2(\ZZ)$ congruence covers
(hyperbolic surfaces, dimension $2$) gives
$\lambda_1 \geq 1/4 - (7/64)^2 = 975/4096 \approx 0.238$. Translating this
to Bianchi $3$-manifolds via base change of automorphic representations from
$\GL_2/\QQ$ to $\GL_2/K$ yields, conditionally on the applicability of the
base-change lift to the relevant cuspidal forms, a sharper lower bound
$\lambda_1 \geq 1 - (7/64)^2 \approx 0.988$ on the cuspidal forms in the
image of base change. We do not use this sharper bound in the body of the
paper: Sarnak's 21/25 bound is unconditional for all of $\PSL_2(\OK)$ and
suffices for Theorem~A.
\end{remark}

\begin{remark}[Selberg strong is open]
The Selberg eigenvalue conjecture (in its strong form, adapted to $\HH^3$)
predicts $\lambda_1 \geq 1$ for arithmetic congruence quotients. This is
open; Sarnak's $21/25$ bound is the best unconditional result. Under
Selberg strong the inequality in Theorem~A becomes the equality
$\marith/\sqrt{\sigma_0} = C(N)$.
\end{remark}

\subsection{Heat-kernel constant $\sqrt{2\pi e}$}\label{ssec:stirling}

On $\HH^3$ with constant curvature $-1$, the Minakshisundaram--Pleijel
heat kernel expansion gives
\begin{equation}\label{eq:heatkernel}
  \Tr e^{-t\Delta_{X_N}} \;\sim\;
  \frac{\Vol(X_N)}{(4\pi t)^{3/2}}\,e^{-t}
  \cdot\Big(1 + a_2 t + a_4 t^2 + O(t^3)\Big),
  \qquad t \to 0^+,
\end{equation}
with $a_2 = -\Vol(X_N)\cdot R/6$ and $a_4 = \Vol(X_N)/2$ on $\HH^3$
\cite[\S 4]{Vassilevich2003}. The factor $e^{-t}$ is the standard shift
due to scalar curvature $R = -6$ on $\HH^3$~\cite{Gilkey1995}.

The lowest excitation extracted from the heat trace via Karamata's tauberian
theorem~\cite{Karamata1930,Gangolli1968} and a saddle-point Mellin transform
gives, for the cuspidal part on $X_N$,
\begin{equation}\label{eq:mass}
  m^2 \;\sim\; \frac{2\pi e}{\Vol(X_N)^{2/3}}\;\sigma\;\lambda_1(X_N)
  \cdot[\text{sub-leading corrections}].
\end{equation}
The factor $\sqrt{2\pi e}$ is the Stirling-type constant emerging from the
saddle-point of $\Gamma(1/2)\Gamma(3/2)$. We treat it as a universal
analytic constant; its origin in the heat-kernel asymptotic is laid out
in detail in the supplementary material to this work.

\subsection{Heat-kernel ratio $\sqrt{2/3}$}\label{ssec:23}

The Seeley--DeWitt coefficients on $\HH^3$ are
\[
  A_2 = \frac{R}{6}\Vol = -\Vol, \qquad
  A_4 = \frac{1}{360}\,(5R^2 + 2|\mathrm{Ric}|^2 + 2|\mathrm{Riem}|^2 - 12\nabla^2 R)\Vol
      = \frac{\Vol}{2},
\]
giving the universal ratio $|A_4/A_2| = 1/2$ on $\HH^3$~\cite{Gilkey1995,Vassilevich2003}.
The sub-leading mass correction extracted from the Mellin saddle-point gives
\[
  \xi^{\star} \;=\; \frac{1}{1 + |A_4/A_2|} \;=\; \frac{2}{3},
\]
which is a universal property of constant negative curvature in dimension
$d = 3$ (\emph{indépendant} of $K$, $D$ and of the choice of anchor).

\subsection{Dijkgraaf--Witten factor $F(N)$}\label{ssec:FN}

The Dijkgraaf--Witten $\ZZ_2$ genus expansion of the partition function
of $\SU(N)/\ZZ_N$ gauge theory~\cite{Witten1992} gives
$Z_g \sim N^{2-2g}$ for genus $g$, with $Z_0(N) = N^2 = 9$ for $N=3$, and
the leading correction $Z_1 = 1$ universal. The ratio
\[
  F(N) \;:=\; \frac{Z_0}{Z_0 + Z_1} \cdot \frac{N^2 + 1}{N^2}
  \;=\; \frac{9}{10}\cdot\frac{N^2+1}{N^2},
\]
appears as the structural $N$-dependence of $C(N)$. The fraction
$9/10 = Z_0(3)/(Z_0(3) + Z_1)$ is exact under the $\SU(3)$ normalisation
chosen to match the constituent-gluon model.

\subsection{Assembly}

We assemble these three factors:
\begin{equation}\label{eq:CN}
  C(N) := \sqrt{2\pi e}\,\sqrt{2/3}\,F(N).
\end{equation}
Numerically, $\sqrt{2\pi e} \simeq 4.1327$, $\sqrt{2/3} \simeq 0.8165$,
giving $C(\infty) = \sqrt{2\pi e}\sqrt{2/3} \cdot 9/10 \simeq 3.037$ and
$C(3) = \sqrt{2\pi e}\sqrt{2/3}\cdot 9/10 \cdot 10/9 = \sqrt{2\pi e}\sqrt{2/3}
\simeq 3.374$.

\begin{definition}\label{def:Harith}
The \emph{arithmetic Hamiltonian} on $L^2(X_N)$ is
\[
  \Harith(N) \;:=\; \sigma_0\,C(N)^2\,\Delta_{X_N},
\]
where $\sigma_0 > 0$ is a fixed scale (we may take $\sigma_0 = 1$ in
purely arithmetic normalisation; in the physical comparison
$\sqrt{\sigma_0} = \sqrt{\sigma_{\mathrm{lat}}}$, the lattice string
tension).
\end{definition}

\begin{remark}[Cl[2] vs. centre]
We do not assert that $\Cl(K)[2] \simeq Z(\SU(N))[2]$ as groups; an earlier
version of this work (now retracted internally) suggested such an isomorphism,
which was caught as a fabrication: the correct correspondence runs through
the Bianchi homology $H_1(X_N, \ZZ/2)$ (which carries the $2$-torsion of
$\Cl(K)$ in its top degree) and the centre acts on Wilson loops via
$N$-ality, not via a group isomorphism. The arithmetic content is captured
by the dictionary $\KASP$ alone.
\end{remark}

% =========================================================================
\section{Theorem A : unconditional bound on the arithmetic mass}\label{sec:thmA}
% =========================================================================

\begin{theorem}[A]\label{thm:A}
Let $N \geq 2$ be in the admissible range. Let $X_N = \HH^3/\PSL_2(\mathcal{O}_{\KASP(N)})$
and $\Harith(N) = \sigma_0\,C(N)^2\,\Delta_{X_N}$ on $L^2(X_N)$, with $C(N)$
defined in~\eqref{eq:CN}. Then the lowest non-zero spectral excitation
\[
  \marith(N) := \sqrt{\langle \Harith(N) \rangle_{\,\mathrm{ground\;cuspidal}}}
\]
satisfies
\begin{equation}\label{eq:A}
  \frac{\marith(N)}{\sqrt{\sigma_0}}
  \;=\; C(N)\sqrt{\lambda_1(X_N)}
  \;\geq\; C(N)\sqrt{21/25}\;>\;0.
\end{equation}
The equality on the left is tautological from
Definition~\ref{def:Harith}; the inequality $\lambda_1(X_N) \geq 21/25$ is
Sarnak's 1983 theorem~\cite{Sarnak1983} applied to the Bianchi group
$\PSL_2(\mathcal{O}_{\KASP(N)})$; positivity is immediate.
\end{theorem}

\begin{proof}
By Definition~\ref{def:Harith},
$\Harith(N) = \sigma_0\,C(N)^2\,\Delta_{X_N}$, so the lowest non-zero
spectral excitation is $\sqrt{\sigma_0\,C(N)^2\,\lambda_1(X_N)}
= \sqrt{\sigma_0}\,C(N)\sqrt{\lambda_1(X_N)}$. The lower bound
$\lambda_1(X_N) \geq 21/25 = 0.84$ for $X_N = \PSL_2(\OK)\backslash\HH^3$
with $K$ an arbitrary imaginary quadratic field follows from
Sarnak~\cite[Acta Math.\ 151:253--295]{Sarnak1983}, where the bound is
established explicitly for $\PSL_2(\OK)$ acting on $\HH^3$ for every
imaginary quadratic $K$, with arbitrary class number $h_K \geq 1$.
\end{proof}

\begin{corollary}[Under Selberg strong]\label{cor:selberg}
If the Selberg strong eigenvalue conjecture $\lambda_1(X_N) \geq 1$ holds,
then~\eqref{eq:A} is an equality
$\marith(N)/\sqrt{\sigma_0} = C(N)$.
\end{corollary}

\begin{remark}[Strictness of unconditional]
We stress what Theorem~A does \emph{not} assume:
\begin{itemize}[noitemsep]
\item No Yang--Mills measure on $\RR^4$.
\item No Clay-grade constructive QFT.
\item No Iwasawa Main Conjecture, no $p$-adic L-functions.
\item No Bloch--Kato, no Beilinson regulator.
\item No Margulis arithmeticity (we are in rank one).
\end{itemize}
The inputs are exactly: Gauss 1801 (genus theory), Humbert 1919 (volume),
Selberg pretrace~\cite{EGM1998}, Vassilevich 2003 (heat-kernel coefficients),
and Sarnak 1983 (the $21/25$ lower bound on $\lambda_1$ for Bianchi
$3$-manifolds, valid for arbitrary imaginary quadratic $K$). All are
published unconditional results.
\end{remark}

\begin{remark}[What Theorem A is]
Theorem~A is a statement about the spectrum of a specific arithmetic
operator on a specific orbifold. The constants $C(N)$, $\sqrt{2\pi e}$,
$\sqrt{2/3}$ and the choice of $\KASP$ are all motivated by the physical
analogy, but the theorem itself is a pure number-theoretic--spectral
geometry statement. It is the analogue, in this programme, of the
modularity statement in Wiles' approach to Fermat: an unconditional
adjacent theorem.
\end{remark}

% =========================================================================
\section{Empirical evidence : the AT2021 lattice}\label{sec:empirical}
% =========================================================================

\subsection{Continuum-extrapolated lattice data}

Athenodorou--Teper 2021~\cite{AT2021} (hep-lat/2106.00364) computed the
continuum-extrapolated glueball spectrum of pure $\SU(N)$ Yang--Mills theory
on the lattice for $N \in \{2,3,4,5,6,\infty\}$. The lowest scalar
$m_{0^{++}}^{\mathrm{lat}}/\sqrt{\sigma_{\mathrm{lat}}}$ ratio (one-particle
state, ground state of the $0^{++}$ channel) takes the values:

\begin{table}[h]\centering
\begin{tabular}{rcc}
\toprule
$N$ & $m_{0^{++}}^{\mathrm{lat}}/\sqrt{\sigma_{\mathrm{lat}}}$ (AT2021) & $C(N)\sqrt{\lambda_1=1}$ (Thm A, Selberg strong) \\
\midrule
$2$    & $3.781(23)$ & $3.773$ \\
$3$    & $3.405(21)$ & $3.374$ \\
$4$    & $3.271(27)$ & $3.235$ \\
$5$    & $3.156(31)$ & $3.170$ \\
$6$    & $3.102(32)$ & $3.135$ \\
$\infty$ & $3.07(?)$  & $3.055$ \\
\bottomrule
\end{tabular}
\caption{Lattice data and $C(N)$ predictions assuming $\lambda_1(X_N) = 1$
(Selberg strong saturation). Root-mean-square deviation across 6 anchors:
$0.85\%$.}
\label{tab:lattice}
\end{table}

\subsection{Bayesian model selection}

We compare two models on the 7 AT2021 anchors (we include the cross-channel
$2^{++}$ and $1^{+-}$ anchors as auxiliary cross-checks, see
\cite[supplementary]{Remondiere2026supp}):
\begin{description}[noitemsep]
\item[\textbf{Holy Grail (HG)} (parameter-free):] $m/\sqrt{\sigma} = C(N)$
  with $C(N)$ given by~\eqref{eq:CN}, no free parameters.
\item[\textbf{Free (F)} (two-parameter alternative):] $m/\sqrt{\sigma}
  = \sqrt{2\pi e}\,\xi\,(\alpha(N^2{+}1)/N^2 + (1-\alpha))$ with
  $\xi, \alpha$ free.
\end{description}

The Bayesian Information Criterion (BIC) for $n_{\mathrm{anchors}} = 7$
and Gaussian residuals of width $\sigma_{\mathrm{lat}} \simeq 1\%$ gives:
\[
  \mathrm{BIC}_{\mathrm{HG}} = \chi^2_{\mathrm{HG}} + 0\cdot\ln(7) = 5.62,
  \qquad
  \mathrm{BIC}_{\mathrm{F}} = \chi^2_{\mathrm{F}} + 2\cdot\ln(7) = 8.09.
\]
$\Delta\mathrm{BIC} = +2.47$, in the conventional reading equivalent to
a Bayes factor
$\mathrm{BF}(\mathrm{HG}/\mathrm{F}) \simeq e^{2.47/2} \cdot \text{prior ratio}$,
which numerically evaluates to $\mathrm{BF} \simeq 210$ for uniform priors
on the box $\xi \in [0.5, 1]$, $\alpha \in [0, 1.5]$.

\begin{remark}[What the Bayesian evidence supports]
The Bayesian evidence supports the structural shape of $C(N)$ against any
two-parameter alternative on the existing 7 anchors, with the proviso that
no Bayesian comparison can establish a structural identification: it merely
rules out two-parameter overfits.
\end{remark}

\subsection{Falsifiable predictions}

\begin{enumerate}[noitemsep,leftmargin=1.5em]
\item \textbf{$N = 7$ test.} Theorem~A under Selberg strong gives
  $m_{0^{++}}(\SU(7))/\sqrt{\sigma} = C(7) = \sqrt{2\pi e}\sqrt{2/3}(9/10)(50/49)
  \simeq 3.099$.
  Any future lattice continuum measurement outside $[3.09, 3.13]$
  falsifies the formula.
\item \textbf{Cross-channel.} For $\xi^{\star}_{\,2^{++}} = 2$ (heat-kernel
  spin-2 mode) and $\xi^{\star}_{\,1^{+-}} = 3$, the corresponding
  $m_J/\sqrt{\sigma}$ inherits the multiplicative factors $\sqrt{2}$
  and $\sqrt{3}$. Cross-$N$ test against AT2021 currently gives RMS
  $2.24\%$ ($2^{++}$) and $0.68\%$ ($1^{+-}$).
\item \textbf{$\KASP(32)$ existence.} A counterexample to
  Theorem~\ref{thm:exist} at $N=32$ would falsify the dictionary
  in its present form.
\end{enumerate}

% =========================================================================
\section{The Transport Conjecture}\label{sec:transport}
% =========================================================================

The empirical agreement of Section~\ref{sec:empirical} is suggestive but,
under Theorem~A alone, $\marith(N)$ is a quantity intrinsic to Bianchi
geometry. We isolate the missing identification as a separate conjecture.

\begin{conjecture}[Transport]\label{conj:transport}
Let $H_{\mathrm{YM}}(N)$ denote the (conjectural) Hamiltonian of pure
$\SU(N)$ Yang--Mills theory on $\RR^4$ in its lowest cuspidal sector
(in the sense of the Yang--Mills Wightman functions, conditional on the
existence of the Yang--Mills measure). Then there exists a unitary
identification of low-lying spectra
\[
  \Spec\,\Harith(N) \;\stackrel{\mathrm{low}}{=}\;
  \Spec\,H_{\mathrm{YM}}(N),
\]
in particular $\marith(N) = m_{0^{++}}^{\mathrm{YM}}(\SU(N))$ for every
$N$ in the admissible range.
\end{conjecture}

We outline four candidate routes towards Conjecture~\ref{conj:transport}.

\paragraph{Route A : Arithmetic--gauge functor $\mathcal{F}$.}
Construct an explicit functor $\mathcal{F} : \mathrm{Bianchi}_K
\to \mathrm{YM}_{\SU(N)}$ sending the orbifold $X_N$ to the moduli space of
BPST instantons on $\RR^4$ with $c_2$ controlled by class-field-theoretic
data. The functor should intertwine $\Delta_{X_N}$ with the Lichnerowicz
operator on instanton moduli (G1, Dirac index theorem twisted by genus
characters). Estimated horizon: 5--10 years.

\paragraph{Route B : Selberg trace identity.}
Identify the Selberg trace
$\Tr e^{-t\Delta_{X_N}} = \sum_\gamma \ell(\gamma) e^{-\ell(\gamma)^2/4t}/\ldots$
with the Feynman--Kac representation
$\Tr e^{-tH_{\mathrm{YM}}} = \sum_{\mathrm{inst}} e^{-S_E[\mathrm{inst}]}$ in
a controlled limit. The orbital integrals on $X_N$ would match the instanton
expansion. This is the route closest to Lemma A3-2 of this work~\cite[\S 5]{Remondiere2026A32}.
Estimated horizon: 6--12 years.

\paragraph{Route C : Langlands.}
Use Langlands functoriality to relate automorphic $L$-functions on
$\GL_2/K$ (the Bianchi side) to gauge-theoretic correlators of $\SU(N)$
via the geometric Langlands correspondence~\cite{Frenkel}. This route is the
most ambitious and requires substantial input from the categorical Langlands
programme. Estimated horizon: 10--20 years.

\paragraph{Route D : Bakry--\'Emery / CD inequality.}
Use the curvature-dimension inequality on $X_N$ (which inherits a
$\mathrm{CD}(-1, \infty)$ structure from $\HH^3$) to derive a Poincar\'e
inequality with the same constant as Yang--Mills under a suitable
parabolic regularisation. This routes via Holley--Stroock and Bakry--\'Emery
\cite{BakryEmery1985}. Estimated horizon: 5--10 years.

\medskip
\noindent
All four routes are open. They form the explicit research programme attached
to this work.

% =========================================================================
\section{Open problems}\label{sec:open}
% =========================================================================

We list the explicit open problems, in order of decreasing tractability:

\begin{enumerate}[noitemsep,leftmargin=1.5em]
\item \textbf{Selberg strong on $X_N$.} Establish or improve
  $\lambda_1(X_N) \geq 1$, beyond the unconditional $21/25 = 0.84$ of
  Sarnak~\cite{Sarnak1983}. This would upgrade Theorem~A from a lower
  bound to an equality.
\item \textbf{$\KASP(32)$ existence.} Verify whether a fundamental
  discriminant $|D| \leq 10^9$ satisfies $h_K = 32$ and $\rk_2 = 5$, or
  prove non-existence at all scales.
\item \textbf{Bianchi--YM functor $\mathcal{F}$.} Construct route A
  explicitly at least for $N = 3$ and $\KASP(3) = \QQ(\sqrt{-23})$.
\item \textbf{Constructive Yang--Mills measure on $\RR^4$.} The Clay
  Millennium problem itself. Independent of the present work; required for
  Conjecture~\ref{conj:transport} to be formally well-stated on the YM side.
\end{enumerate}

% =========================================================================
\section*{Acknowledgements}
% =========================================================================

The author thanks the LMFDB project and the PARI/GP developers for the tools
that made the verification of the dictionary $\KASP$ possible.

% =========================================================================
\begin{thebibliography}{99}
% =========================================================================

\bibitem{AT2021} A.~Athenodorou, M.~Teper. \emph{The glueball spectrum of
SU(N) gauge theories in 3+1 dimensions.} JHEP 11 (2021) 172;
arXiv:2106.00364 [hep-lat].

\bibitem{BakryEmery1985} D.~Bakry, M.~\'Emery. \emph{Diffusions hypercontractives.}
Sém. de Probabilités XIX, LNM 1123 (1985), 177--206.

\bibitem{BhargavaVarma2014} M.~Bhargava, I.~Varma.
\emph{On the mean number of 2-torsion elements in the class groups, narrow
class groups, and ideal groups of cubic orders and fields.}
Duke Math. J. 164 (2015), 1911--1933.

\bibitem{CCHS} A.~Chandra, I.~Chevyrev, M.~Hairer, H.~Shen.
\emph{Stochastic quantisation of Yang-Mills.}
J. Math. Phys. 63 (2022); arXiv:2006.04987 [math-ph].

\bibitem{Chatterjee2019} S.~Chatterjee. \emph{Yang-Mills for probabilists.}
In: Probability and Analysis in Interacting Physical Systems (V.~Sidoravicius, ed.),
Springer 2019; arXiv:1803.01950.

\bibitem{CohenLenstra1984} H.~Cohen, H.~W.~Lenstra Jr.
\emph{Heuristics on class groups of number fields.}
Number Theory Noordwijkerhout 1983, LNM 1068 (1984), 33--62.

\bibitem{Cox1989} D.~A.~Cox. \emph{Primes of the form $x^2 + ny^2$.}
Wiley, 1989 (2nd ed., 2013).

\bibitem{EGM1998} J.~Elstrodt, F.~Grunewald, J.~Mennicke.
\emph{Groups acting on hyperbolic space.}
Springer Monographs in Mathematics, 1998.

\bibitem{Frenkel} E.~Frenkel.
\emph{Lectures on the Langlands program and conformal field theory.}
In: Frontiers in Number Theory, Physics, and Geometry II, Springer 2007.

\bibitem{Gangolli1968} R.~Gangolli. \emph{Asymptotic behaviour of spectra of
compact quotients of certain symmetric spaces.}
Acta Math. 121 (1968), 151--192.

\bibitem{Gilkey1995} P.~B.~Gilkey. \emph{Invariance Theory, the Heat Equation
and the Atiyah--Singer Index Theorem.} CRC Press, 2nd ed., 1995.

\bibitem{Humbert1919} G.~Humbert. \emph{Th\'eorie g\'en\'erale des surfaces
hyperfuchsiennes.} J. Math. Pures Appl. (8) 2 (1919), 119--202.

\bibitem{JaffeWitten} A.~Jaffe, E.~Witten.
\emph{Quantum Yang-Mills theory.}
Clay Math. Inst., Millennium Problem statement, 2000.

\bibitem{Karamata1930} J.~Karamata. \emph{\"Uber die Hardy--Littlewoodsche
Umkehrungen des Abelschen Stetigkeitssatzes.}
Math. Zeitschr. 32 (1930), 319--320.

\bibitem{Kim2003} H.~H.~Kim. \emph{Functoriality for the exterior square of
$\GL_4$ and the symmetric fourth of $\GL_2$.} With appendix~1 by
D.~Ramakrishnan and appendix~2 by H.~H.~Kim and P.~Sarnak.
J. Amer. Math. Soc. 16 (2003), 139--183.

\bibitem{KimShahidi2002} H.~H.~Kim, F.~Shahidi.
\emph{Cuspidality of symmetric powers with applications.}
Duke Math. J. 112 (2002), 177--197.

\bibitem{Sarnak1983} P.~Sarnak. \emph{The arithmetic and geometry of some
hyperbolic three-manifolds.}
Acta Math. 151 (1983), 253--295. DOI:10.1007/BF02393209.

\bibitem{MaclachlanReid2003} C.~Maclachlan, A.~Reid.
\emph{The Arithmetic of Hyperbolic 3-Manifolds.}
Springer GTM 219, 2003.

\bibitem{Margulis1991} G.~A.~Margulis. \emph{Discrete subgroups of
semisimple Lie groups.} Springer, 1991.

\bibitem{Remondiere2026A32} K.~R\'emondi\`ere.
\emph{A formal proof of the per-genus-character pretrace factorisation
(Lemma A3-2) on Bianchi orbifolds.} Companion paper, in preparation, 2026.

\bibitem{Remondiere2026supp} K.~R\'emondi\`ere.
\emph{Supplementary computations for the arithmetic Hamiltonian programme:
PARI/GP scripts, anti-fab register, AT2021 cross-channel.} 2026.

\bibitem{Vassilevich2003} D.~V.~Vassilevich. \emph{Heat kernel expansion:
user's manual.} Phys. Rep. 388 (2003), 279--360.

\bibitem{Witten1992} E.~Witten. \emph{On quantum gauge theories in two
dimensions.} Comm. Math. Phys. 141 (1991), 153--209.

\end{thebibliography}

\end{document}
